% Options for packages loaded elsewhere
% Options for packages loaded elsewhere
\PassOptionsToPackage{unicode}{hyperref}
\PassOptionsToPackage{hyphens}{url}
\PassOptionsToPackage{dvipsnames,svgnames,x11names}{xcolor}
%
\documentclass[
  letterpaper,
  DIV=11,
  numbers=noendperiod]{scrartcl}
\usepackage{xcolor}
\usepackage{amsmath,amssymb}
\setcounter{secnumdepth}{5}
\usepackage{iftex}
\ifPDFTeX
  \usepackage[T1]{fontenc}
  \usepackage[utf8]{inputenc}
  \usepackage{textcomp} % provide euro and other symbols
\else % if luatex or xetex
  \usepackage{unicode-math} % this also loads fontspec
  \defaultfontfeatures{Scale=MatchLowercase}
  \defaultfontfeatures[\rmfamily]{Ligatures=TeX,Scale=1}
\fi
\usepackage{lmodern}
\ifPDFTeX\else
  % xetex/luatex font selection
\fi
% Use upquote if available, for straight quotes in verbatim environments
\IfFileExists{upquote.sty}{\usepackage{upquote}}{}
\IfFileExists{microtype.sty}{% use microtype if available
  \usepackage[]{microtype}
  \UseMicrotypeSet[protrusion]{basicmath} % disable protrusion for tt fonts
}{}
\makeatletter
\@ifundefined{KOMAClassName}{% if non-KOMA class
  \IfFileExists{parskip.sty}{%
    \usepackage{parskip}
  }{% else
    \setlength{\parindent}{0pt}
    \setlength{\parskip}{6pt plus 2pt minus 1pt}}
}{% if KOMA class
  \KOMAoptions{parskip=half}}
\makeatother
% Make \paragraph and \subparagraph free-standing
\makeatletter
\ifx\paragraph\undefined\else
  \let\oldparagraph\paragraph
  \renewcommand{\paragraph}{
    \@ifstar
      \xxxParagraphStar
      \xxxParagraphNoStar
  }
  \newcommand{\xxxParagraphStar}[1]{\oldparagraph*{#1}\mbox{}}
  \newcommand{\xxxParagraphNoStar}[1]{\oldparagraph{#1}\mbox{}}
\fi
\ifx\subparagraph\undefined\else
  \let\oldsubparagraph\subparagraph
  \renewcommand{\subparagraph}{
    \@ifstar
      \xxxSubParagraphStar
      \xxxSubParagraphNoStar
  }
  \newcommand{\xxxSubParagraphStar}[1]{\oldsubparagraph*{#1}\mbox{}}
  \newcommand{\xxxSubParagraphNoStar}[1]{\oldsubparagraph{#1}\mbox{}}
\fi
\makeatother


\usepackage{longtable,booktabs,array}
\usepackage{calc} % for calculating minipage widths
% Correct order of tables after \paragraph or \subparagraph
\usepackage{etoolbox}
\makeatletter
\patchcmd\longtable{\par}{\if@noskipsec\mbox{}\fi\par}{}{}
\makeatother
% Allow footnotes in longtable head/foot
\IfFileExists{footnotehyper.sty}{\usepackage{footnotehyper}}{\usepackage{footnote}}
\makesavenoteenv{longtable}
\usepackage{graphicx}
\makeatletter
\newsavebox\pandoc@box
\newcommand*\pandocbounded[1]{% scales image to fit in text height/width
  \sbox\pandoc@box{#1}%
  \Gscale@div\@tempa{\textheight}{\dimexpr\ht\pandoc@box+\dp\pandoc@box\relax}%
  \Gscale@div\@tempb{\linewidth}{\wd\pandoc@box}%
  \ifdim\@tempb\p@<\@tempa\p@\let\@tempa\@tempb\fi% select the smaller of both
  \ifdim\@tempa\p@<\p@\scalebox{\@tempa}{\usebox\pandoc@box}%
  \else\usebox{\pandoc@box}%
  \fi%
}
% Set default figure placement to htbp
\def\fps@figure{htbp}
\makeatother





\setlength{\emergencystretch}{3em} % prevent overfull lines

\providecommand{\tightlist}{%
  \setlength{\itemsep}{0pt}\setlength{\parskip}{0pt}}



 


\KOMAoption{captions}{tableheading}
\makeatletter
\@ifpackageloaded{caption}{}{\usepackage{caption}}
\AtBeginDocument{%
\ifdefined\contentsname
  \renewcommand*\contentsname{Table of contents}
\else
  \newcommand\contentsname{Table of contents}
\fi
\ifdefined\listfigurename
  \renewcommand*\listfigurename{List of Figures}
\else
  \newcommand\listfigurename{List of Figures}
\fi
\ifdefined\listtablename
  \renewcommand*\listtablename{List of Tables}
\else
  \newcommand\listtablename{List of Tables}
\fi
\ifdefined\figurename
  \renewcommand*\figurename{Figure}
\else
  \newcommand\figurename{Figure}
\fi
\ifdefined\tablename
  \renewcommand*\tablename{Table}
\else
  \newcommand\tablename{Table}
\fi
}
\@ifpackageloaded{float}{}{\usepackage{float}}
\floatstyle{ruled}
\@ifundefined{c@chapter}{\newfloat{codelisting}{h}{lop}}{\newfloat{codelisting}{h}{lop}[chapter]}
\floatname{codelisting}{Listing}
\newcommand*\listoflistings{\listof{codelisting}{List of Listings}}
\makeatother
\makeatletter
\makeatother
\makeatletter
\@ifpackageloaded{caption}{}{\usepackage{caption}}
\@ifpackageloaded{subcaption}{}{\usepackage{subcaption}}
\makeatother
\usepackage{bookmark}
\IfFileExists{xurl.sty}{\usepackage{xurl}}{} % add URL line breaks if available
\urlstyle{same}
\hypersetup{
  pdftitle={Course Topics},
  colorlinks=true,
  linkcolor={blue},
  filecolor={Maroon},
  citecolor={Blue},
  urlcolor={Blue},
  pdfcreator={LaTeX via pandoc}}


\title{Course Topics}
\author{}
\date{}
\begin{document}
\maketitle

\renewcommand*\contentsname{Table of contents}
{
\hypersetup{linkcolor=}
\setcounter{tocdepth}{3}
\tableofcontents
}

\section{BUS 411 --- Fixed Income}\label{bus-411-fixed-income}

\subsection{Lecture Topics}\label{lecture-topics}

This document lists the planned lecture topics and subtopics for the
course.

\subsection{Lecture 1: History and Role of Fixed Income
Markets}\label{lecture-1-history-and-role-of-fixed-income-markets}

\subsubsection{1.1 What is fixed income?}\label{what-is-fixed-income}

\begin{itemize}
\tightlist
\item
  Definition of fixed income securities
\item
  Debt versus equity financing
\item
  Main issuers and investors
\end{itemize}

\subsubsection{1.2 Historical development of fixed income
markets}\label{historical-development-of-fixed-income-markets}

\begin{itemize}
\tightlist
\item
  Early sovereign borrowing
\item
  Development of government debt markets
\item
  Growth of corporate bond markets
\end{itemize}

\subsubsection{1.3 Economic role of fixed income
securities}\label{economic-role-of-fixed-income-securities}

\begin{itemize}
\tightlist
\item
  Funding governments
\item
  Funding corporations
\item
  Saving, retirement, and liability management
\end{itemize}

\begin{center}\rule{0.5\linewidth}{0.5pt}\end{center}

\subsection{Lecture 2: Time Value of Money and
Discounting}\label{lecture-2-time-value-of-money-and-discounting}

\subsubsection{2.1 Time value of money}\label{time-value-of-money}

\begin{itemize}
\tightlist
\item
  Present value
\item
  Future value
\item
  Discount factors
\end{itemize}

\subsubsection{2.2 Compounding
conventions}\label{compounding-conventions}

\begin{itemize}
\tightlist
\item
  Annual compounding
\item
  Semiannual compounding
\item
  Continuous compounding
\end{itemize}

\subsubsection{2.3 Discounting fixed income cash
flows}\label{discounting-fixed-income-cash-flows}

\begin{itemize}
\tightlist
\item
  Cash flow timelines
\item
  Pricing as present value of future cash flows
\item
  Relation between discount rates and value
\end{itemize}

\begin{center}\rule{0.5\linewidth}{0.5pt}\end{center}

\subsection{Lecture 3: Annuities and Liability
Valuation}\label{lecture-3-annuities-and-liability-valuation}

\subsubsection{3.1 Fixed annuities}\label{fixed-annuities}

\begin{itemize}
\tightlist
\item
  Ordinary annuities
\item
  Annuities due
\item
  Perpetuities
\end{itemize}

\subsubsection{3.2 Variable annuities}\label{variable-annuities}

\begin{itemize}
\tightlist
\item
  Basic structure
\item
  Investment-linked cash flows
\item
  Valuation considerations
\end{itemize}

\subsubsection{3.3 Defined benefit pension
plans}\label{defined-benefit-pension-plans}

\begin{itemize}
\tightlist
\item
  Nature of pension liabilities
\item
  Discounting promised payments
\item
  Funding and valuation issues
\end{itemize}

\subsubsection{3.4 Liability valuation}\label{liability-valuation}

\begin{itemize}
\tightlist
\item
  Present value of liabilities
\item
  Matching assets to liabilities
\item
  Liability sensitivity to interest rates
\end{itemize}

\begin{center}\rule{0.5\linewidth}{0.5pt}\end{center}

\subsection{Lecture 4: Government
Bonds}\label{lecture-4-government-bonds}

\subsubsection{4.1 Types of government
securities}\label{types-of-government-securities}

\begin{itemize}
\tightlist
\item
  Treasury bills
\item
  Treasury notes
\item
  Treasury bonds
\end{itemize}

\subsubsection{4.2 Features of government
bonds}\label{features-of-government-bonds}

\begin{itemize}
\tightlist
\item
  Face value
\item
  Coupon rate
\item
  Maturity
\item
  Clean price and dirty price
\end{itemize}

\subsubsection{4.3 Valuation of government
bonds}\label{valuation-of-government-bonds}

\begin{itemize}
\tightlist
\item
  Pricing from discounted cash flows
\item
  Yield to maturity
\item
  Accrued interest
\end{itemize}

\subsubsection{4.4 Government bond
markets}\label{government-bond-markets}

\begin{itemize}
\tightlist
\item
  Primary market
\item
  Secondary market
\item
  Benchmark role of sovereign debt
\end{itemize}

\begin{center}\rule{0.5\linewidth}{0.5pt}\end{center}

\subsection{Lecture 5: Corporate Bonds and Credit
Risk}\label{lecture-5-corporate-bonds-and-credit-risk}

\subsubsection{5.1 Corporate bond
features}\label{corporate-bond-features}

\begin{itemize}
\tightlist
\item
  Secured and unsecured bonds
\item
  Seniority
\item
  Covenants
\end{itemize}

\subsubsection{5.2 Valuation of corporate
bonds}\label{valuation-of-corporate-bonds}

\begin{itemize}
\tightlist
\item
  Pricing with promised cash flows
\item
  Credit spread over government rates
\item
  Yield spread interpretation
\end{itemize}

\subsubsection{5.3 Credit risk}\label{credit-risk}

\begin{itemize}
\tightlist
\item
  Default risk
\item
  Recovery risk
\item
  Credit ratings
\end{itemize}

\subsubsection{5.4 Corporate bond market
considerations}\label{corporate-bond-market-considerations}

\begin{itemize}
\tightlist
\item
  Liquidity
\item
  Issuer quality
\item
  Term to maturity and risk
\end{itemize}

\begin{center}\rule{0.5\linewidth}{0.5pt}\end{center}

\subsection{Lecture 6: Yield Measures and the Term Structure of Interest
Rates}\label{lecture-6-yield-measures-and-the-term-structure-of-interest-rates}

\subsubsection{6.1 Yield measures}\label{yield-measures}

\begin{itemize}
\tightlist
\item
  Current yield
\item
  Yield to maturity
\item
  Holding period return
\end{itemize}

\subsubsection{6.2 Spot and forward rates}\label{spot-and-forward-rates}

\begin{itemize}
\tightlist
\item
  Spot rates
\item
  Forward rates
\item
  Discount function
\end{itemize}

\subsubsection{6.3 Yield curve construction and
interpretation}\label{yield-curve-construction-and-interpretation}

\begin{itemize}
\tightlist
\item
  Normal, flat, and inverted curves
\item
  Economic interpretation
\item
  Expectations about future rates
\end{itemize}

\subsubsection{6.4 Term structure
concepts}\label{term-structure-concepts}

\begin{itemize}
\tightlist
\item
  Expectations view
\item
  Liquidity preference
\item
  Market segmentation intuition
\end{itemize}

\begin{center}\rule{0.5\linewidth}{0.5pt}\end{center}

\subsection{Lecture 7: Duration}\label{lecture-7-duration}

\subsubsection{7.1 Interest rate risk in fixed
income}\label{interest-rate-risk-in-fixed-income}

\begin{itemize}
\tightlist
\item
  Bond price and yield relationship
\item
  Why prices fall when yields rise
\end{itemize}

\subsubsection{7.2 Macaulay duration}\label{macaulay-duration}

\begin{itemize}
\tightlist
\item
  Definition
\item
  Weighted average time to cash flow receipt
\item
  Interpretation
\end{itemize}

\subsubsection{7.3 Modified duration}\label{modified-duration}

\begin{itemize}
\tightlist
\item
  Definition
\item
  Approximate price sensitivity
\item
  Duration as first-order risk measure
\end{itemize}

\subsubsection{7.4 Determinants of
duration}\label{determinants-of-duration}

\begin{itemize}
\tightlist
\item
  Coupon
\item
  Maturity
\item
  Yield level
\end{itemize}

\begin{center}\rule{0.5\linewidth}{0.5pt}\end{center}

\subsection{Lecture 8: Convexity}\label{lecture-8-convexity}

\subsubsection{8.1 Limits of duration-only
approximation}\label{limits-of-duration-only-approximation}

\begin{itemize}
\tightlist
\item
  Nonlinearity in bond prices
\item
  Curvature of the price-yield relationship
\end{itemize}

\subsubsection{8.2 Convexity measure}\label{convexity-measure}

\begin{itemize}
\tightlist
\item
  Definition
\item
  Interpretation
\item
  Why convexity matters
\end{itemize}

\subsubsection{8.3 Duration-convexity
approximation}\label{duration-convexity-approximation}

\begin{itemize}
\tightlist
\item
  First-order and second-order effects
\item
  Improved bond price change estimates
\item
  Application to larger yield changes
\end{itemize}

\subsubsection{8.4 Comparing bonds using
convexity}\label{comparing-bonds-using-convexity}

\begin{itemize}
\tightlist
\item
  Price sensitivity differences
\item
  Trade-off between yield and convexity
\item
  Portfolio implications
\end{itemize}

\begin{center}\rule{0.5\linewidth}{0.5pt}\end{center}

\subsection{Lecture 9: Fixed Income Risk
Management}\label{lecture-9-fixed-income-risk-management}

\subsubsection{9.1 Duration-based risk
management}\label{duration-based-risk-management}

\begin{itemize}
\tightlist
\item
  Measuring portfolio interest rate exposure
\item
  Portfolio duration
\end{itemize}

\subsubsection{9.2 Immunization}\label{immunization}

\begin{itemize}
\tightlist
\item
  Asset-liability matching
\item
  Duration matching
\item
  Protecting against small parallel shifts in yields
\end{itemize}

\subsubsection{9.3 Convexity in risk
management}\label{convexity-in-risk-management}

\begin{itemize}
\tightlist
\item
  Why duration matching alone is not enough
\item
  Role of convexity in immunization
\item
  Rebalancing issues
\end{itemize}

\subsubsection{9.4 Practical limitations}\label{practical-limitations}

\begin{itemize}
\tightlist
\item
  Non-parallel yield curve shifts
\item
  Model risk
\item
  Transaction costs
\end{itemize}

\begin{center}\rule{0.5\linewidth}{0.5pt}\end{center}

\subsection{Lecture 10: Bonds with
Contingencies}\label{lecture-10-bonds-with-contingencies}

\subsubsection{10.1 Bonds with embedded
options}\label{bonds-with-embedded-options}

\begin{itemize}
\tightlist
\item
  Callable bonds
\item
  Putable bonds
\item
  Convertible features as a general idea
\end{itemize}

\subsubsection{10.2 Valuation issues for contingent
bonds}\label{valuation-issues-for-contingent-bonds}

\begin{itemize}
\tightlist
\item
  Cash flows depend on future events
\item
  Option-like features in bond valuation
\item
  Impact on price and yield
\end{itemize}

\subsubsection{10.3 Mortgage-backed and prepayable
securities}\label{mortgage-backed-and-prepayable-securities}

\begin{itemize}
\tightlist
\item
  Prepayment risk
\item
  Negative convexity
\item
  Valuation challenges
\end{itemize}

\subsubsection{10.4 Risk characteristics}\label{risk-characteristics}

\begin{itemize}
\tightlist
\item
  Reinvestment risk
\item
  Call risk
\item
  Extension and contraction risk
\end{itemize}

\begin{center}\rule{0.5\linewidth}{0.5pt}\end{center}

\subsection{Lecture 11: Fixed Income Derivative
Securities}\label{lecture-11-fixed-income-derivative-securities}

\subsubsection{11.1 Overview of fixed income
derivatives}\label{overview-of-fixed-income-derivatives}

\begin{itemize}
\tightlist
\item
  Why derivatives are used
\item
  Hedging, speculation, and portfolio management
\end{itemize}

\subsubsection{11.2 Interest rate forwards and
futures}\label{interest-rate-forwards-and-futures}

\begin{itemize}
\tightlist
\item
  Basic contract structure
\item
  Hedging interest rate exposure
\item
  Relation to underlying fixed income securities
\end{itemize}

\subsubsection{11.3 Interest rate swaps}\label{interest-rate-swaps}

\begin{itemize}
\tightlist
\item
  Fixed-for-floating exchange
\item
  Basic valuation intuition
\item
  Uses in risk management
\end{itemize}

\subsubsection{11.4 Options on interest
rates}\label{options-on-interest-rates}

\begin{itemize}
\tightlist
\item
  Caps
\item
  Floors
\item
  Basic role in managing rate risk
\end{itemize}

\begin{center}\rule{0.5\linewidth}{0.5pt}\end{center}

\subsection{Coverage of Core Course
Themes}\label{coverage-of-core-course-themes}

The course is organized around three major themes:

\subsubsection{1. History of fixed income
securities}\label{history-of-fixed-income-securities}

\begin{itemize}
\tightlist
\item
  Covered primarily in Lecture 1
\end{itemize}

\subsubsection{2. Valuation methods for fixed income
securities}\label{valuation-methods-for-fixed-income-securities}

\begin{itemize}
\tightlist
\item
  Fixed and variable annuities: Lecture 3
\item
  Government bonds: Lecture 4
\item
  Corporate bonds: Lecture 5
\item
  Defined benefit pension plans: Lecture 3
\item
  Bonds with contingencies: Lecture 10
\item
  Fixed income derivative securities: Lecture 11
\end{itemize}

\subsubsection{3. Fixed income risk management using duration and
convexity}\label{fixed-income-risk-management-using-duration-and-convexity}

\begin{itemize}
\tightlist
\item
  Duration: Lecture 7
\item
  Convexity: Lecture 8
\item
  Risk management and immunization: Lecture 9
\end{itemize}




\end{document}
